% =====================================================================
%  conclusions.tex
%  Revised 13 September 2026 for the inclusion of Campaign 9.
%  Labels defined: ch:conclusions, sec:concl-answer, sec:concl-reform,
%  sec:concl-candidates, sec:concl-limits, sec:concl-future,
%  eq:concl-reform, tab:concl-future.
%  Style: short paragraphs, no semicolons, no em dashes, formulas alone
%  with term lists, enumerations as lists.
% =====================================================================

\chapter{Conclusions}
\label{ch:conclusions}

\section{Summary of the study and its findings}
\label{sec:concl-answer}

The study asked which core designs a LABGENE-class pressurised water
reactor of \SI{48}{MW} thermal could adopt when the only information
available is the open record. It answered by building a
surrogate-assisted multi-objective optimisation around a Monte Carlo
model of the core and running it through nine campaigns, 636
evaluations in all. Each campaign changed one element of the method and
measured its effect. Campaigns 1 to 8 maximised the cycle length and
minimised the radial peaking factor. The eighth produced a front of
eight designs and a second set of four designs controllable with two
regulating banks, and the difference between the two sets led to the
reformulated problem of Campaign 9.

Four findings carry the weight.

\begin{enumerate}
  \item A derived bound is never a substitute for a measured
        constraint. The reactivity ceiling of Campaigns 1 to 6 was
        replaced in Campaign 7 by a directly measured controllability
        screen, and the comparison showed the derived ceiling
        rejecting fourteen designs whose own measured bank worth made
        them controllable (Section~\ref{sec:res-c7-screens}). Campaign
        8 reversed the nesting so that the measurement does the work
        and the bound is inert. The retrospective derivation also gave
        the historic ceiling of 1.35 its provenance, as the all-bank
        controllability bound of this core to within two per cent.

  \item The finite-height correction is a constant of the geometry.
        The axial leakage factor varies by \SI{160}{pcm} over a factor
        of two in excess reactivity, so a single value corrects the
        whole design space and every cycle length of Campaigns 8 and 9
        refers to a finite-height core (Section~\ref{sec:meth-lax}).

  \item The cycle length of this core is a requirement to be met. Every
        Campaign 8 design exceeded the reported refuelling interval,
        and the excess reactivity spent on the longer cycle had to be
        held by the control rods (Section~\ref{sec:concl-reform}).

  \item At the required cycle length, the boron the core needs falls
        with its gadolinia loading. Among the feasible Campaign 9
        designs between 4.15 and \SI{4.55}{wt\%}, the critical
        concentration at the most reactive state falls from 3003 to
        \SI{1406}{ppm} as the gadolinia rises from 0.14 to
        \SI{4.42}{wt\%} (Section~\ref{sec:res-c9-gd}), a trade that the
        objectives of the first eight campaigns could not reward.
\end{enumerate}

\section{Reformulation of the cycle length as a constraint}
\label{sec:concl-reform}

The optimisation was posed, from the first campaign, as the
maximisation of cycle length against the minimisation of the radial
peaking factor. That is the natural formulation for a naval core,
whose defining advantage is time between refuellings, and it was kept
through eight campaigns so that the campaigns remain comparable.
Campaign 8 showed that it is the wrong formulation for the final design
choice, for a reason that only became visible once controllability was
measured rather than assumed.

The open record gives the prototype a refuelling interval of five
years~\cite{dinizcosta_2017}. Through Equation~\eqref{eq:mission-efpd} that interval
corresponds to 913 to 1461 effective full power days for capacity
factors between 0.5 and 0.8, and no open source states the capacity
factor. Every design on the Campaign 8 front exceeds the upper end of
that range, the shortest by a third and the longest by a factor of
four (Section~\ref{sec:res-c8-mission}). The objective kept pushing
cycle length long after the requirement was met, and it could only do
so by raising the beginning-of-life excess reactivity, because a
longer cycle is bought with reactivity to be burned later.

That excess has to be held by something at the start of life. In this
model the soluble boron and the gadolinia are already present in every
eigenvalue, so what remains is held by the control rods. Every design
on the front needs the four regulating banks to hold between
\num{8930} and \SI{13627}{pcm} at power, which means banks inserted
deep into the core throughout early life. A pressurised water reactor
does not operate that way. Its regulating banks sit near the withdrawn
position at power and carry only the power defect and the manoeuvring
band, while the excess is held by boron adjustment and by burnable
absorber whose burnout tracks the fuel. The four Campaign 8 designs
that behave that way, controllable by two banks and holding \num{2564}
to \SI{4068}{pcm}, still deliver 1142 to 3381 effective full power
days (Section~\ref{sec:res-c8-split}).

The requirement on cycle length is a floor. A cycle beyond it gives the
mission nothing and costs excess reactivity that the control system
must absorb. The peaking factor and the reactivity to be held at
start-up have no such floor, and both improve in the direction of the
design goal for as long as the search can move them. Those are the
objectives, and the cycle length belongs among the constraints, at the
value the mission requires. Campaign 9 was posed that way.

\begin{equation}
\begin{aligned}
\min_{\mathbf{x}}\ &\big(F_{\Delta H}(\mathbf{x}),\ c_\mathrm{max}(\mathbf{x})\big) \\
\text{subject to}\ &\mathrm{EFPD}(\mathbf{x}) \ge \mathrm{EFPD}_\mathrm{req},\quad
g_j(\mathbf{x}) \le 0
\end{aligned}
\label{eq:concl-reform}
\end{equation}

\noindent where the symbols have the following meaning.

\begin{itemize}
  \item $F_{\Delta H}$ is the radial enthalpy-rise hot channel factor
        with all rods out, dimensionless, the first objective.

  \item $c_\mathrm{max}$ is the soluble boron concentration at which
        the core is critical at the most reactive state of its cycle,
        in ppm, the second objective, defined in
        Equation~\eqref{eq:c9-cmax}.

  \item $\mathrm{EFPD}_\mathrm{req}$ is the required cycle length,
        1826 effective full power days, five years at a capacity factor
        of one, the conservative choice of Section~\ref{sec:meth-c9}.

  \item $g_j$ are the remaining constraints of
        Section~\ref{sec:meth-c9}: the reactivity window, the
        enrichment cap, the vessel fit, the peaking screen at 1.65, and
        the control screen at beginning of life and at the operating
        maximum.
\end{itemize}

The second objective carries the reactivity that the Campaign 8 front
could not hold. The critical concentration at the operating maximum
expresses the beginning-of-life excess and the gadolinium hump in one
quantity, and it does so in the unit in which the moderator
temperature coefficient bounds it: a design whose critical
concentration exceeds the boron ceiling of
Section~\ref{sec:res-c8-mtc} cannot be held at power by boron with a
non-positive coefficient.

Scored on the Campaign 8 archive without transport, the reformulated
problem at the five-year floor selects designs 47 and 13
(Section~\ref{sec:res-c8-retro}). Campaign 9 then searched the same
design space with the boron measured on every design. Its front holds
five designs at peaking factors of 1.490 to 1.545 and boron demands of
1406 to \SI{1908}{ppm}, every one within 121 days of the floor, so the
constraint binds (Section~\ref{sec:res-c9-front}). At that floor the
enrichment of the front is 4.2 to \SI{4.5}{wt\%}, below the 6 to
\SI{8}{wt\%} stated for the prototype in 2014 and far below the
\SI{19}{wt\%} reported in 2025.

\section{Candidate designs}
\label{sec:concl-candidates}

The two formulations give two stages of candidate.

Under the formulation of Campaigns 1 to 8, the candidate is design 47
of Table~\ref{tab:c8-front}, the low-peaking floor of the front
(Section~\ref{sec:res-c8-decision}). Its cycle length is 1941\,EFPD,
reproduced within 1.5 per cent by a replicate depletion, and its
radial peaking factor is 1.510, with a seed-to-seed standard deviation
of about 0.014 at the archive fidelity that does not separate it from
its three nearest neighbours on that axis. On the three-dimensional
core model with structural components at 1000\,ppm it holds a two-bank margin of
1330\,pcm, which the gadolinium hump leaves intact because beginning of
life is its most reactive state. Its differential boron worth is
6.98\,pcm per ppm, its rods-only margin under the four regulating
banks with no boron is 3069\,pcm, and its critical concentration of
\SI{2327}{ppm} lies below the boron ceiling at both primary conditions
(Sections~\ref{sec:res-c8-boron}, \ref{sec:res-c8-zeroboron}
and~\ref{sec:res-c8-mtc}). Design 53 holds more margin and needs less
boron but meets five years only at a capacity factor of 0.8, and design
31 loses its two-bank margin at the maximum of its hump.

On the objectives of Campaign 9, both Campaign 8 designs selected by
the retrospective are dominated by members of the Campaign 9 front.
Design 47 of Campaign 8 needs 419 to \SI{602}{ppm} more boron than
designs 35, 40 and 34 of Campaign 9, at a peaking factor that differs
by less than one and a half seed standard deviations, and design 13 is
dominated by designs 47 and 44 of Campaign 9 on both objectives
(Section~\ref{sec:res-c9-compare}). The comparison uses critical
concentrations obtained by different routes and is indicative until the
two archives are scored on the same footing.

Within Campaign 9 the champion is C9-47, the lowest boron demand of the
front, and C9-47 and C9-34 are the only members that stay on the front
in every resample of the measurement noise
(Section~\ref{sec:res-c9-convergence}). Two further designs, C9-29 and
C9-30, are excluded from the nominal front by less than the noise, so the
candidate set of the campaign holds seven designs. Four of the five front
members, C9-47 among them, would have been buried two to four fronts deep
under the cycle-length objectives (Section~\ref{sec:res-c9-discarded}).

Re-solved with all rods out on the three-dimensional model, C9-47 and
C9-34 stay on the front. C9-40 enters it, because its peaking factor
falls by 0.040, and C9-44 leaves it, because C9-47 then has both the lower
peaking and the lower boron demand. Both changes exceed four standard
deviations of the seed noise, but the one that brings C9-40 in rests on
the design with the largest seed spread of the five
(Section~\ref{sec:res-c9-peaking}). Neither C9-47 nor C9-34 is yet a final
candidate. Three of the checks that the Campaign 8 designs passed are now passed by
the Campaign 9 front. On the three-dimensional model with structural
components, at each design's own critical
concentration, the first two regulating banks hold every front member
subcritical with 9228 to \SI{9545}{pcm} of margin
(Section~\ref{sec:res-c9-front}), and the boron ceiling measured on the
lattice of each candidate is cleared, by \SI{1491}{ppm} for C9-47 and
\SI{1027}{ppm} for C9-34 (Section~\ref{sec:res-c9-ceiling}). The k-target
leakage factor of every front design is constant over the cycle within its
uncertainty (Section~\ref{sec:res-c9-post}). Until the remaining checks of
that section are closed, design 47 of Campaign 8 is the most completely confirmed
design of the study, and the Campaign 9 front is the set from which the
final choice will be made.

\section{Limitations}
\label{sec:concl-limits}

\begin{enumerate}
  \item The model is a two-dimensional beginning-of-life core in the
        optimisation loop, corrected for axial leakage by a measured
        constant. On the eleven Campaign 8 designs the
        three-dimensional core model gives the peaking, the rod worth
        and the axial leakage with rods withdrawn and inserted, and
        shows that the leakage correction of the unrodded core does not
        hold for the rodded core. On the five Campaign 9 front designs
        it gives the peaking, the axial power shape, the bank worth and
        the axial rod worth curves. No depletion was run in three
        dimensions in any campaign.

  \item The control screen tests hold-down at power with the operating
        boron present and no xenon. It is not a shutdown-margin
        demonstration, which needs the cold state, all assemblies
        minus the highest-worth one, and a stated margin.

  \item Up to Campaign 8 the reactivity swing over the cycle was not
        penalised. The optimiser removed burnable absorber wherever it
        could, and humps of several thousand pcm appear in the depletion
        histories. Campaign 9 scores the boron demand at the operating
        maximum and applies the control screen there, which closes this
        limitation for the last campaign only.

  \item Every cycle length comes from a depletion at a constant
        \SI{1000}{ppm} and is a lower bound for a core operated with
        boron letdown.

  \item The statistical error on the peaking objective was not
        measured during the campaigns, only bounded. Measured afterwards
        at the confirmation statistics, it is 0.011 on the Campaign 8
        designs and 0.0061 on the Campaign 9 front, and the ordering of
        neighbouring front designs on that axis is not resolved in
        either campaign. The peaking objective is also two-dimensional,
        and on the Campaign 9 front the three-dimensional value changes
        which designs are non-dominated.

  \item Every design below a pitch of \SI{1.2242}{cm} in Campaigns 1
        to 7 carried clipped guide tubes. Campaigns 8 and 9 are free of
        the defect and the earlier results are caveated.

  \item Thermal-hydraulic feedback is absent. Temperatures are fixed at
        representative hot full power values, so no departure from
        nucleate boiling ratio and no fuel temperature limit is
        evaluated. The moderator temperature coefficient was evaluated
        outside the loop, on one Campaign 8 design and on eight
        Campaign 9 lattices, and enters only through the boron ceiling.

  \item The boron ceiling was measured on the lattice of Campaign 8
        design 47, with \SI{2.88}{wt\%} gadolinia on twelve rods. The
        Campaign 9 front carries 2.6 to \SI{4.4}{wt\%} on 16 or 20
        rods, so the ceiling is a reference for that front until it is
        measured on its own lattice.
\end{enumerate}

\section{Continuation of this work}
\label{sec:concl-future}

The work stops where the deadline put it, with the post-analysis of
Campaign 9 in progress. Table~\ref{tab:concl-future} lists the
continuations in the order they would be taken up, each with what it
needs and what it would settle.

\begin{table}[htbp]
  \centering
  \footnotesize
  \caption{Continuations of the study, in order of priority.}
  \label{tab:concl-future}
  \begin{tabularx}{\textwidth}{@{}>{\raggedright\arraybackslash}p{3.6cm}>{\raggedright\arraybackslash}X>{\raggedright\arraybackslash}p{4.0cm}@{}}
    \toprule
    Continuation & What it involves & What it settles \\
    \midrule
    Campaign 9 post-analysis
      & Boron ceiling at 15.5\,MPa on C9-34 and C9-44, independent boron worth
      & The final candidate \\
    \addlinespace
    Shutdown margin
      & Cold zero-power state, all 32 assemblies minus the highest-worth one, operating boron, margin of at least about \SI{1600}{pcm}
      & Whether the candidates meet a licensing-style shutdown criterion rather than an operating one \\
    \addlinespace
    Boron letdown in the depletion
      & A critical boron search at each burnup step in place of the constant \SI{1000}{ppm}
      & Cycle lengths of an operated core in place of lower bounds \\
    \addlinespace
    Core-level depletion
      & Depletion of the core in place of the reflected assembly, at the cost estimated in Section~\ref{sec:meth-dep-level}
      & The radial burnup distribution, and the removal of the leakage-corrected target approximation \\
    \addlinespace
    Burnable absorber design
      & Mixed gadolinia concentrations per assembly, enriched gadolinium, or erbia, following the soluble-boron-free literature cited in the review
      & A flatter swing at the same cycle length, and the path to a boron-free variant \\
    \addlinespace
    Peaking uncertainty in the loop
      & Read the core tally errors on every evaluation and carry them into the surrogate as noise
      & Statistically resolved ordering of the front on the peaking axis \\
    \addlinespace
    Acquisition robustness
      & A per-variable minimum spread in the batch selection, so that a batch cannot satisfy the separation rule by moving two variables
      & Ends the collapse mode seen in Campaigns 6 and 7 \\
    \addlinespace
    Thermal-hydraulic coupling
      & OpenMC coupled to a subchannel or porous-medium solver, in the frameworks surveyed in the review
      & Coupled temperature feedback, departure from nucleate boiling, and fuel temperature margins on the candidates \\
    \addlinespace
    Multi-batch loading
      & Reload strategy and loading-pattern optimisation for a core that refuels rather than replaces
      & The peaking gap between a single-batch uniform core and licensed practice \\
    \bottomrule
  \end{tabularx}
\end{table}

The first continuation is scripted and was running when this chapter
was written. Its last stage is the gate for any further campaign. If
the ratio of three- to two-dimensional peaking is constant across the
Campaign 9 archive, the three-dimensional objective is a monotone
transform of the two-dimensional one, the front does not move and no
further campaign is needed. On the eleven Campaign 8 designs the ratio
was not constant, with a Spearman coefficient of 0.909 and a largest
swap of three positions. On the five Campaign 9 front designs the ratio
is 0.992 with a standard deviation of 0.014, and substituting it already
changes the non-dominated set of those five, so the measurement on all
60 designs is what decides whether a three-dimensional campaign is
needed.

The study set out to show that a core can be designed from four open
facts and a declared design space, and that a surrogate-assisted loop
can carry the search at a cost a workstation affords. Both were shown
within the limits stated above: a two-dimensional beginning-of-life
model in the loop, a fresh single-batch core, fixed temperatures and a
constant boron concentration in the depletion. Within those limits the
framework produced a front of buildable, controllable designs, the
finding that the cycle length of this core is a requirement to be met,
and a reformulated campaign whose front meets that requirement with
less boron than the two Campaign 8 designs selected by the same
formulation. What lies beyond the limits, the three-dimensional power
shape of the new front, the shutdown margin, the thermal-hydraulic
feedback and the reload strategy, is the work of
Table~\ref{tab:concl-future}, and the archive, the checkpoints and the
scripts that would carry it out are in place.
